We are going to prove the above statement using weak induction. For n >= 1, let P(n) be the statement

n(n + 1) / 2

Basis Step: We will now prove the base case P(1)

P(1) = 1(1 + 1)/2 = 1

Since G is a simple undirected graph, G with 1 node has 0 edges. Since 1 is the upper bound and 0 <= 1, P(1) holds.

Induction Step: We will now prove the inductive case

For each k >= 1, assume the claim is true for P(k) and prove that it is ture for P(k + 1). The induction hypothesis states that

A simple graph with k + 1 nodes has at most (k + 1)(k + 2)/2 edges

So, take some graph G with k + 1 nodes. Choose one node v, and remove it, and all edges connected to v. We are left with a simple graph containing k edges. By our induction, the graph has at most

k(k + 1)/2 edges

Now put v back. Because the graph is simple v can't have any edges to itself or multiple edges. Since there are k other edges, v can add at most k edges.

Simplify (k(k + 1)/2) + k

(k(k + 1)/2) + k = (k(k + 1) + 2k)/2
= (k^2 + 3k)/2

Our desired upper bound is (k + 1)(k + 2)/2 = (k^2 + 3k + 2)/2

Clearly (k^2 + 3k)/2 <= (k^2 + 3k + 2)/2

Hence, E(G) <= (k^2 + 3k + 2)/2, and P(k + 1) holds

Therefore, by weak mathematical induction, P(n) is true for all n >= 1
